%%% FORMATO E IDIOMA %%%
    \documentclass[letterpaper,DIV=15,12pt]{scrartcl}
    \usepackage[spanish,mexico,shorthands=off,es-lcroman]{babel}
    \usepackage{scrlayer-scrpage}
        \clearpairofpagestyles
        \ihead{\footnotesize \textit{Teoría de Conjuntos III}}
        \ohead{\footnotesize \textit{2026-2}}
        \cfoot{\normalfont\thepage}
        \addtokomafont{title}{\bfseries \rmfamily}
        \setlength{\headsep}{5pt}
        \newpairofpagestyles{beginstyle}{
            \clearpairofpagestyles
            \KOMAoptions{headsepline=false}
            \cfoot{\footnotesize \pagemark}
            \cfoot{\normalfont\thepage}
        }
        \addtokomafont{section}{\raggedleft\rmfamily}
        \addtokomafont{subsection}{\raggedleft\rmfamily}
        \addtokomafont{subsubsection}{\raggedleft\rmfamily}
        \setlength{\headsep}{8pt}
        \setlength{\footskip}{20pt}
        \setlength{\textheight}{600pt}
%%% PAQUETES %%%
    \usepackage{ifthen}
    \usepackage[shortlabels]{enumitem}
        \setenumerate[1]{label=\MakeLowercase{\roman*}), ref=\roman*}
        \setenumerate[2]{label=\MakeLowercase{\alph*}), ref=\alph*}
    \usepackage{amsmath}
    \usepackage{amsthm}\usepackage[T1]{fontenc}
    \usepackage{stix2}
    \usepackage{multicol}
	\renewcommand{\qedsymbol}{$\blacksquare$}
	\addto\captionsspanish{\renewcommand{\proofname}{\normalfont\bfseries Demostración}}	
	\newenvironment{solucion}{\renewcommand{\qedsymbol}{$\boxtimes$}\begin{proof}[\normalfont\bfseries Solución]}{\end{proof}}
%%% E J E R C I C I O S %%%
    \newcounter{Ejer}
    \newcounter{Pts}
    \setcounter{Ejer}{1}
    \setcounter{Pts}{1}
    \newcommand{\pts}{}
    \newenvironment{ejercicio}[1]{\noindent
        \ifthenelse{\equal{#1}{1}}{\renewcommand{\pts}{\textbf{(#1 pt)}}}{\renewcommand{\pts}{\textbf{(#1 pts)}}}\textbf{Ej. \theEjer} \pts\stepcounter{Ejer}}{}
%%% C O M A N D O S %%%
    \renewcommand{\emptyset}{\varnothing}
    \newcommand{\tq}{:}
%% D O C U M E N T O %%
\begin{document}
\thispagestyle{plain}

    \begin{center}
        {\fontsize{30}{60}\rmfamily \textbf{Fórmulas}} \\ \vspace{.2cm} Teoría de Conjuntos III, 2026-2
    \end{center}
    \begin{flushright}
        \footnotesize \hfill Profesor: Luis Jesús Turcio Cuevas.\\
        \hfill Ayudante: Hugo Víctor García Martínez.
    \end{flushright}
        
    La siguiente lista de fórmulas son (equivalentes a una fórmula) $\Delta_0$, y por tanto, son absolutas para clases transitivas.
    \begin{multicols}{2}
        \begin{enumerate}
            \item $x=0$.
            \item $0 \in x$.
            \item $x=\{a,b\}$.
            \item $x=(a,b)$.
            \item $x \subseteq y$.
            \item $x$ es transitivo.
            \item $x$ es ordinal.
            \item $x$ es límite.
            \item $x$ es natural.
            \item $x$ es inductivo.
            \item $x \subseteq \omega$.
            \item $x = \omega$.
            \item $x = \bigcup a$.
            \item $x = \bigcap a$.
            \item $x = a \cup \{a\}$.
            \item $x = a \cup b$.
            \item $x = a \cap b$.
            \item $x = a \setminus b$.
            \item $x = a \times b$.
            \item $R$ es relación.
            \item $R$ es relación transitiva.
            \item $R$ es relación simétrica.
            \item $R$ es relación reflexiva.
            \item $R$ es relación asimétrica.
            \item $R$ es relación tricotómica.
            \item $R$ es orden parcial en $P$.
            \item $R$ es orden parcial irreflexivo en $P$.
            \item $R$ es orden parcial en $P$ y $x R y$.
            \item $R$ es orden parcial irreflexivo en $P$ y $x R y$.
            \item $R$ es orden total en $P$.
            \item $R$ es orden total en $P$ y $x R y$.
            \item $f$ es función, denotado $\operatorname*{fun}(f)$.
            \item $\operatorname*{fun}(f)$ y $y=\operatorname*{dom}(f)$.
            \item $\operatorname*{fun}(f)$ y $y=\operatorname*{ima}(f)$.
            \item $\operatorname*{fun}(f)$ y $y=f(x)$.
            \item $\operatorname*{fun}(f)$ y $g=f \upharpoonright A$.
            \item $\operatorname*{fun}(f)$ y $f$ es sobreyeciva.
            \item $\operatorname*{fun}(f)$ y $f$ es inyectiva.
            \item $\operatorname*{fun}(f)$ y $f$ es biyectiva.
        \end{enumerate}
    \end{multicols}
\end{document}
